\exercisegroup

\begin{enumerate}
\def\labelenumi{\arabic{enumi})}
\item
 设 X 是拓扑空间，\(x, y\) 是 X 中的点。证明映射 \(p_x\) 和 \(p_y\) 分别以 \(x\) 和 \(y\) 为像，且为常值映射；它们同伦，当且仅当 \(x\) 和 \(y\) 属于 X 的同一道路连通分支。
\item
  设 S 是从 \([0,1]\) 到 R 的函数 \(x \mapsto \sin(\pi/x)\) 的图形；设 \(\overline{S}\) 是它在 \(R^{2}\) 中的闭包。
\end{enumerate}

\begin{enumerate}
\def\labelenumi{\alph{enumi})}
\item
  证明空间 S 和 \(\overline{S}\) 是连通的。
\item
  证明 \(\overline{S}\) 中任何包含点 \((1,0)\) 和 \((0,1)\) 的连通紧子集必等于 S；但是，不存在一个从 I 到 \(\overline{S}\) 的同胚把 0 映到 \((1,0)\)，把 1 映到 \((0,1)\)。
\item
  求空间 \(\overline{S}\) 的道路连通分支。
\item
  证明空间 \(\overline{\mathrm{S}}\) 不是局部连通的。
\end{enumerate}

\begin{enumerate}
\def\labelenumi{\arabic{enumi})}
\setcounter{enumi}{2}
\item
  设 \(a \in \overline{\mathrm{S}}\)。证明，带基点空间 \((\overline{\mathrm{S}}, a)\) 有万有覆叠，当且仅当 \(a \in \{0\} \times [-1, 1]\)。
\item
  令 \(D = [0, 1] \cap (\mathbf{R} - \mathbf{Q})\)。
\end{enumerate}

\begin{enumerate}
\def\labelenumi{\alph{enumi})}
\item
 令 \(\varphi\colon D \to [0,1]\) 为一个映射，令 X 为 \([0,1]^2\) 中由所有点 \((x,\varphi(x))\)（其中 \(x\in\mathrm{D}\)）构成的集合的补集。证明 X 是连通的。
\item
  证明存在双射 \(\Phi\)，其值域是所有起点为 \((0,0)\)、终点为 \((1,1)\) 且位于 \([0,1]^2\) 中的道路所成的集合，定义域为 D。对所有 \(x \in \mathrm{D}\)，令 \(\psi(x)\) 为横坐标为 \(x\) 的点，该点位于道路 \(\Phi(x)\) 上。令 Y 为 \([0,1]^2\) 中由所有点 \(\psi(x)\)（其中 \(x \in \mathrm{D}\)）构成的集合的补集。证明空间 Y 连通且局部连通，但不是道路连通的。
\end{enumerate}

\begin{enumerate}
\def\labelenumi{\arabic{enumi})}
\setcounter{enumi}{4}
\tightlist
\item
  设 r 是 \(N \cup \{\infty, \omega\}\) 中的元素，n 和 d 是整数，X 是 \(C^{r}\) 类的 \(R^{n}\) 中的子流形，且是闭的、纯维数为 \(d\)。\(^{3}\)
\end{enumerate}

\begin{enumerate}
\def\labelenumi{\alph{enumi})}
\item
  证明存在 X 在 \(\mathbf{R}^n\) 中的一个邻域 U，以及一个 \(\mathrm{C}^r\) 类同构，从 \(\mathrm{X} \times ] - 1; 1[^{n - d}\) 到 \(\mathrm{U}\)。
\item
  X 中的每条道路都严格同伦于一条 \(C^{r}\) 类道路。
\item
  若 X 中两条 \(C^{r}\) 类道路同伦，则它们可由一个 \(C^{r}\) 类同伦连接。
\end{enumerate}

\begin{enumerate}
\def\labelenumi{\arabic{enumi})}
\setcounter{enumi}{5}
\item
  设 U 是数值空间 \(R^{n}\) 的一个开子集。若道路 \(c: I \to U\) 存在有限序列 \((t_{0}, \ldots, t_{n})\)，满足 \(0 = t_{0} < t_{1} < \cdots < t_{n} = 1\)，且 c 在每个区间 \([t_{i-1}, t_{i}]\) 上的限制都是仿射的，则称道路 c 是分段仿射的。证明 U 中的每条道路都严格同伦于一条分段仿射道路。
\item
  \begin{enumerate}
  \def\labelenumii{\alph{enumii})}
  \tightlist
  \item
    设 R 是一个集合。若 S 是 R 的子集，我们用 \(j_{S}\) 表示从 \(I^{S}\) 到 \(I^{R}\) 的映射，它把函数 \(f: S \to I\) 映为从 R 到 I 的唯一函数：该函数在 S 上的限制等于 f，并且在 R-S 的每一点都为零。
  \end{enumerate}
\end{enumerate}

证明 \(j_{S}\) 是从 \(I^{S}\) 到 \(I^{R}\) 中由在 S 外为零的函数组成的子空间上的同胚。

\begin{enumerate}
\def\labelenumi{\alph{enumi})}
\setcounter{enumi}{1}
\tightlist
\item
  令 X 为从 R 到 I 的、在 R 的某个可数子集之外为零的映射所成的集合。我们把空间族 \(I^{S}\) 的归纳极限记为 Y，其中 S 遍历 R 的所有可数子集，并在 Y 上赋予由乘积拓扑得到的归纳极限拓扑。
\end{enumerate}

从 \(I^{S}\) 到 \(I^{R}\) 的典范映射诱导出从 Y 到 X 的连续双射 j。若 R 不可数，则此双射不是同胚。

\begin{enumerate}
\def\labelenumi{\alph{enumi})}
\setcounter{enumi}{2}
\tightlist
\item
  证明，对每条道路 \(c: I \to X\)，都存在 R 的一个可数子集 S，使得 \(c(t)(s) = 0\) 对所有 \(t \in I\) 及每个 \(s \in R - S\) 都成立。
\end{enumerate}

\(d)\) 推出映射 \(j^{-1}\) 在弧的意义下是连续的。

\begin{enumerate}
\def\labelenumi{\arabic{enumi})}
\setcounter{enumi}{7}
\tightlist
\item
  设 X 是基数 \(\geqslant 2\) 的连通紧拓扑空间。假设对 X 中每一对不同点 \((x,y)\)，空间 \(X - \{x,y\}\) 都不连通。
\end{enumerate}

\begin{enumerate}
\def\labelenumi{\alph{enumi})}
\item
  证明对每个 \(x \in \mathrm{X}\)，空间 \(\mathrm{X} - \{x\}\) 是连通的。
\item
  设 a、b 是 X 中的不同点，U、V 是 X 的非空开闭子集，且 \(X - \{a,b\}=U\cup V\)。证明 \(\overline{U} = U \cup \{a, b\}\)、\(\overline{V} = V \cup \{a, b\}\)，并证明这两个集合都是连通的。
\item
  假设 X 有一个可数稠密子集。证明 X 同胚于圆 \(S_{1}\)。（证明存在从 I 到 \(\overline{U}\) 的同胚，把 0 映到 a，把 1 映到 b。）
\end{enumerate}

\begin{enumerate}
\def\labelenumi{\arabic{enumi})}
\setcounter{enumi}{8}
\tightlist
\item
  设 A 是一个不可数良序集，并且其中每个区间 \([u, v]\) 都是可数的。
\end{enumerate}

\begin{enumerate}
\def\labelenumi{\alph{enumi})}
\item
  证明集合 \(L = A \times [0, 1[\)，按字典序全序，并赋予拓扑 \(\mathcal{T}_{0}(L)\)（TG，I，第 91 页，习题 5），是连通的（TG，IV，第 48 页，习题 7）。空间 L 称为 Alexandroff 直线。
\item
  证明对每个点 \(x \in L\)，拓扑空间 \(L - \{x\}\) 都不连通。
\item
  令 \(\overline{L}\) 为在 L 中加入点 \(-\infty\) 和 \(+\infty\) 得到的有序集，使得 \(-\infty < x < +\infty\) 对所有 \(x \in L\) 都成立。赋予它拓扑 \(\mathcal{T}_{0}(\overline{\mathrm{L}})\)；它称为完备 Alexandroff 直线。证明 \(\overline{L}\) 是紧且连通的。
\end{enumerate}

d) 证明不存在从 \([0,1]\) 到 \(\overline{\mathrm{L}}\) 的同胚，把 0 映到 \(-\infty\)，把 1 映到 \(+\infty\)。

\begin{enumerate}
\def\labelenumi{\arabic{enumi})}
\setcounter{enumi}{9}
\tightlist
\item
  令 \(\overline{L}\) 为完备 Alexandroff 直线，令 S 为把 \(\overline{L}\) 中的点 \(-\infty\) 和 \(+\infty\) 识别所得的拓扑空间。
\end{enumerate}

\begin{enumerate}
\def\labelenumi{\alph{enumi})}
\item
  证明空间 S 是紧且连通的，基数至少为 2。
\item
  证明对 S 中每一对不同点 \((x,y)\)，拓扑空间 \(S-\{x,y\}\) 都不连通。
\item
  证明空间 S 不同胚于圆。
\end{enumerate}

\begin{enumerate}
\def\labelenumi{\arabic{enumi})}
\setcounter{enumi}{10}
\tightlist
\item
  \begin{enumerate}
  \def\labelenumii{\alph{enumii})}
  \tightlist
  \item
    令 X 为将空间 \(\mathbf{R} \times \{1, 2\}\) 按满足 \((t, 1)\) 与 \((t, 2)\) 等价（若 \(t \in \mathbf{R}^*\)）的最细等价关系取商所得的空间。证明 X 中不存在一条单射道路，其端点是点 \((0, 1)\) 和 \((0, 2)\) 的像。
  \end{enumerate}
\end{enumerate}

\begin{enumerate}
\def\labelenumi{\alph{enumi})}
\setcounter{enumi}{1}
\tightlist
\item
  令 \(c: I \to S^{1}\) 为由 \(c(t) = e^{3\pi it}\) 定义的道路；它连接点 1 和点 -1。证明 \(S_{1}\) 中不存在严格同伦于道路 c 的单射道路。
\end{enumerate}

\begin{enumerate}
\def\labelenumi{\arabic{enumi})}
\setcounter{enumi}{11}
\item
  设 X 是仿紧空间，并且每一点都有一个同胚于完备度量空间的邻域。证明 X 上存在一个与其拓扑相容的度量，使 X 成为完备度量空间（参见 TG，IX，第 110 页，习题 34）。
\item
  设 X 是分离拓扑空间；假设 X 的每一点都有一个邻域 U，使得 U 中每一对点都能由 U 中的一条单射道路连接。若 X 连通，证明（不使用 III，第 282 页的命题 18）X 是道路连通的，并且 X 中每一对点都能由一条单射道路连接。
\item
  设 X 是集合，c 是无限基数，且 c \textless{} Card(X)。
\end{enumerate}

\begin{enumerate}
\def\labelenumi{\alph{enumi})}
\item
  令 \(\mathcal{U}_{\mathfrak{c}}\) 为 X 的子集 V 所成的集合，其中满足 \(\operatorname{Card}(X - V) \leqslant \mathfrak{c}\) 或 \(V = \varnothing\)。证明 \(\mathcal{U}_{\mathfrak{c}}\) 是 \(X\) 上的一个拓扑。
\item
  证明赋予此拓扑的集合 X 是连通且局部连通的，但其道路连通分支都退化为单点。
\item
  证明锥 C(X) 是道路连通且局部连通的，但不是局部道路连通的。这说明在 III，第 267 页的推论 2 中，不能省略空间局部紧且可度量化这一假设。
\end{enumerate}

